\chapter{1993年京湘鄂云琼贵卷（文）}
\section{单选题}

\begin{problem}
函数 \(f(x)=\sin x+\cos x\) 的最小正周期是（\(\quad\)）

\choices
    {\(2\pi\)}
    {\(2\sqrt{2}\pi\)}
    {\(\pi\)}
    {\(\frac{\pi}{4}\)}
\end{problem}

\begin{answer}
C
\end{answer}

\begin{solution}
利用和角公式，\(f(x)=\sqrt{2}\sin\left(x+\frac{\pi}{4}\right)\). 正弦函数的最小正周期为 \(2\pi\)，平移不改变周期，故最小正周期为 \(2\pi\)，选 C.
\end{solution}

\begin{problem}
如果双曲线的焦距为 \(6\)，两条准线间的距离为 \(4\)，那么该双曲线的离心率为（\(\quad\)）

\choices
    {\(\frac{3}{2}\)}
    {\(\frac{\sqrt{3}}{2}\)}
    {\(\frac{\sqrt{6}}{2}\)}
    {\(2\)}
\end{problem}

\begin{answer}
C
\end{answer}

\begin{solution}
设双曲线的半焦距为 \(c\)，半长轴为 \(a\)，离心率为 \(e\). 由焦距为 \(6\)，得 \(2c=6\)，即 \(c=3\). 双曲线两条准线的距离为 \(2a/e\)，所以 \(2a/e=4\). 又 \(e=c/a\)，于是
\(\frac{2a^2}{c}=4\)，\(\Longrightarrow\)，\(a^2=2c=6\).
因此 \(e=\dfrac{c}{a}=\dfrac{3}{\sqrt{6}}=\dfrac{\sqrt{6}}{2}\)，选 C.
\end{solution}

\begin{problem}
和直线 \(3x-4y+5=0\) 关于 \(x\) 轴对称的直线的方程为（\(\quad\)）

\choices
    {\(3x+4y-5=0\)}
    {\(3x+4y+5=0\)}
    {\(-3x+4y-5=0\)}
    {\(-3x+4y+5=0\)}
\end{problem}

\begin{answer}
B
\end{answer}

\begin{solution}
关于 \(x\) 轴对称把原直线上的点 \((x,y)\) 变为 \((x,-y)\). 在原方程中以 \(-y\) 代替 \(y\)，得
\[
3x-4(-y)+5=0
\]
即 \(3x+4y+5=0\)，选 B.
\end{solution}

\begin{problem}
\(\i^{2n-3}+\i^{2n-1}+\i^{2n+1}+\i^{2n+3}\) 的值为（\(\quad\)）

\choices
    {\(-2\)}
    {\(0\)}
    {\(2\)}
    {\(4\)}
\end{problem}

\begin{answer}
B
\end{answer}

\begin{solution}
相邻两项的指数相差 \(2\)，而 \(\i^{m+2}=-\i^m\). 因而
\(\i^{2n-3}+\i^{2n-1}=0\)，\(\i^{2n+1}+\i^{2n+3}=0\).
四项之和为 \(0\)，选 B.
\end{solution}

\begin{problem}
\(y=x^{\frac{3}{5}}\) 在 \([-1,1]\) 上是（\(\quad\)）

\choices
    {增函数且是奇函数}
    {增函数且是偶函数}
    {减函数且是奇函数}
    {减函数且是偶函数}
\end{problem}

\begin{answer}
A
\end{answer}

\begin{solution}
因为五次方根函数在实数域上严格递增，立方函数也严格递增，所以 \(x^{3/5}=\left(\sqrt[5]{x}\right)^3\) 在 \([-1,1]\) 上为增函数. 又
\[
(-x)^{3/5}=\left(\sqrt[5]{-x}\right)^3=-\left(\sqrt[5]{x}\right)^3=-x^{3/5}
\]
故它是奇函数，选 A.
\end{solution}

\begin{problem}
\(\displaystyle\lim\limits_{n\to\infty}\frac{5n^2-1}{2n^2-n+5}\) 的值为（\(\quad\)）

\choices
    {\(-\frac{1}{5}\)}
    {\(-\frac{5}{2}\)}
    {\(\frac{1}{5}\)}
    {\(\frac{5}{2}\)}
\end{problem}

\begin{answer}
D
\end{answer}

\begin{solution}
分子、分母同除以 \(n^2\)，得
\[
\lim_{n\to\infty}\frac{5-\frac{1}{n^2}}{2-\frac{1}{n}+\frac{5}{n^2}}=\frac{5}{2}
\]
故选 D.
\end{solution}

\begin{problem}
集合 \(M=\left\{x \mid x=\frac{k\pi}{2}+\frac{\pi}{4}, k\in\Z\right\}\)，\(N=\left\{x \mid x=\frac{k\pi}{4}+\frac{\pi}{2}, k\in\Z\right\}\)，则（\(\quad\)）

\choices
    {\(M=N\)}
    {\(M\supseteq N\)}
    {\(M\subseteq N\)}
    {\(M\cap N=\varnothing\)}
\end{problem}

\begin{answer}
C
\end{answer}

\begin{solution}
将 \(N\) 中的通式改写为
\[
x=\frac{k\pi}{4}+\frac{\pi}{2}=\frac{(k+2)\pi}{4}
\]
由于 \(k+2\) 遍历全部整数，\(N=\{m\pi/4\mid m\in\Z\}\). 而 \(M\) 中的系数为 \(2k+1\)，始终为奇数，所以 \(M\subseteq N\). 又 \(0\in N\) 而 \(0\notin M\)，故为真子集，选 C.
\end{solution}

\begin{problem}
\(\sin 20^{\circ}\cos 70^{\circ}+\sin 10^{\circ}\sin 50^{\circ}\) 的值是（\(\quad\)）

\choices
    {\(\frac{1}{4}\)}
    {\(\frac{\sqrt{3}}{2}\)}
    {\(\frac{1}{2}\)}
    {\(\frac{\sqrt{3}}{4}\)}
\end{problem}

\begin{answer}
A
\end{answer}

\begin{solution}
由于 \(\cos70^\circ=\sin20^\circ\)，原式为
\[
\sin^2 20^\circ+\sin10^\circ\sin50^\circ
\]
由积化和差公式，
\(\sin^2 20^\circ=\frac{1-\cos40^\circ}{2}\)，\(\sin10^\circ\sin50^\circ=\frac{\cos40^\circ-\cos60^\circ}{2}\).
相加得 \(\dfrac{1-\frac12}{2}=\dfrac14\)，选 A.
\end{solution}

\begin{problem}
圆 \(x^2+y^2=1\) 上的点到直线 \(3x+4y-25=0\) 的距离的最小值是（\(\quad\)）

\choices
    {\(6\)}
    {\(4\)}
    {\(5\)}
    {\(1\)}
\end{problem}

\begin{answer}
B
\end{answer}

\begin{solution}
圆心为原点，半径为 \(1\). 对圆上点 \((x,y)\)，到直线的距离为
\[
\frac{|3x+4y-25|}{5}
\]
由柯西不等式，\(3x+4y\leq\sqrt{3^2+4^2}\sqrt{x^2+y^2}=5\)，且等号可取. 因此 \(3x+4y-25<0\)，距离的最小值为
\[
\frac{25-5}{5}=4
\]
选 B.
\end{solution}

\begin{problem}
若 \(a\)，\(b\) 是任意实数，且 \(a>b\)，则（\(\quad\)）

\choices
    {\(a^2>b^2\)}
    {\(\frac{b}{a}<1\)}
    {\(\lg \left(a-b\right)>0\)}
    {\(\left(\frac{1}{2}\right)^a<\left(\frac{1}{2}\right)^b\)}
\end{problem}

\begin{answer}
D
\end{answer}

\begin{solution}
逐项检查条件 \(a>b\). 例如取 \(a=0\)，\(b=-1\)，不能推出 \(a^2>b^2\)，故第一项错；第二项在 \(a=0\) 时无意义；\(a-b>0\) 但不一定大于 \(1\)，故第三项错. 因为底数 \(0<\frac12<1\)，指数函数 \(\left(\frac12\right)^x\) 严格递减，所以 \(a>b\) 时
\[
\left(\frac12\right)^a<\left(\frac12\right)^b
\]
故选 D.
\end{solution}

\begin{problem}
一动圆与两圆 \(x^2+y^2=1\) 和 \(x^2+y^2-8x+12=0\) 都外切，则动圆圆心轨迹为（\(\quad\)）

\choices
    {圆}
    {椭圆}
    {双曲线的一支}
    {抛物线}
\end{problem}

\begin{answer}
C
\end{answer}

\begin{solution}
第二个圆可化为
\[
(x-4)^2+y^2=4
\]
其圆心为 \(F(4,0)\)，半径为 \(2\). 设动圆圆心为 \(P\)，半径为 \(r\). 外切条件给出
\(PO=r+1\)，\(PF=r+2\).
相减得 \(PF-PO=1\). 因为两定点 \(F\)，\(O\) 的距离为 \(4\)，所以 \(P\) 的轨迹是以 \(F\)，\(O\) 为焦点、焦点距离之差为定值 \(1\) 的双曲线的一支，选 C.
\end{solution}

\begin{problem}
圆柱轴截面的周长 \(l\) 为定值，那么圆柱体积的最大值是（\(\quad\)）

\choices
    {\(\left(\frac{l}{6}\right)^3\pi\)}
    {\(\frac{1}{9}\left(\frac{l}{2}\right)^3\pi\)}
    {\(\left(\frac{l}{4}\right)^3\pi\)}
    {\(2\left(\frac{l}{4}\right)^3\pi\)}
\end{problem}

\begin{answer}
A
\end{answer}

\begin{solution}
设圆柱底面半径为 \(r\)，高为 \(h\). 轴截面是边长为 \(2r\)，\(h\) 的矩形，故
\(4r+2h=l\)，\(h=\frac l2-2r\)，\(0<r<\frac l4\).
体积
\[
V=\pi r^2h=\pi r^2\left(\frac l2-2r\right)
\]
求导得 \(V'=\pi r(l-6r)\). 当 \(0<r<l/6\) 时，\(V'>0\)；当 \(l/6<r<l/4\) 时，\(V'<0\)，所以 \(V\) 在 \(r=l/6\) 处取得最大值，此时 \(h=l/6\). 因而
\[
V_{\max}=\pi\left(\frac l6\right)^2\frac l6=\left(\frac l6\right)^3\pi
\]
选 A.
\end{solution}

\begin{problem}
\(\left(\sqrt{x}+1\right)^4\left(x-1\right)^5\) 展开式中 \(x^4\) 的系数为（\(\quad\)）

\choices
    {\(-40\)}
    {\(10\)}
    {\(40\)}
    {\(45\)}
\end{problem}

\begin{answer}
D
\end{answer}

\begin{solution}
在 \((\sqrt{x}+1)^4\) 中，要产生整数次幂 \(x^4\)，只能取 \((\sqrt{x})^j\) 中的 \(j=0\)，\(2\)，\(4\). 因而 \(x^4\) 的系数为
\[
\mathrm C_4^0\mathrm C_5^4(-1)^1
 +\mathrm C_4^2\mathrm C_5^3(-1)^2
 +\mathrm C_4^4\mathrm C_5^2(-1)^3
=-5+60-10=45
\]
故选 D.
\end{solution}

\begin{problem}
直角梯形的一个内角为 \(45^{\circ}\)，下底长为上底长的 \(\frac{3}{2}\)，这个梯形绕下底所在的直线旋转一周所成的旋转体的全面积为 \(\left(5+\sqrt{2}\right)\pi\)，则旋转体的体积为（\(\quad\)）

\choices
    {\(2\pi\)}
    {\(\frac{4+\sqrt{2}}{3}\pi\)}
    {\(\frac{5+\sqrt{2}}{3}\pi\)}
    {\(\frac{7}{3}\pi\)}
\end{problem}

\begin{answer}
D
\end{answer}

\begin{solution}
设上底为 \(a\). 因为直角梯形的一个锐角为 \(45^\circ\)，且下底比上底长 \(a/2\)，所以高也为 \(a/2\). 绕下底旋转后得到一个半径为 \(a/2\)、高为 \(a\) 的圆柱和一个底面半径为 \(a/2\)、高为 \(a/2\) 的圆锥.

旋转体全面积由圆柱侧面、圆柱的一个底面和圆锥侧面组成. 斜边长为 \(a/\sqrt{2}\)，故
\[
S=2\pi\frac a2 a+\pi\left(\frac a2\right)^2+\pi\frac a2\frac a{\sqrt2}
=\frac{5+\sqrt2}{4}\pi a^2
\]
由题设 \(S=(5+\sqrt2)\pi\)，得 \(a=2\). 因而旋转体体积为
\[
V=\pi\cdot1^2\cdot2+\frac13\pi\cdot1^2\cdot1=\frac73\pi
\]
选 D.
\end{solution}

\begin{problem}
已知 \(a_1\)，\(a_2\)，\(\cdots\)，\(a_8\) 为各项都大于零的等比数列，公比 \(q\neq 1\)，则（\(\quad\)）

\choices
    {\(a_1+a_8>a_4+a_5\)}
    {\(a_1+a_8<a_4+a_5\)}
    {\(a_1+a_8=a_4+a_5\)}
    {\(a_1+a_8\) 和 \(a_4+a_5\) 的大小关系不能由已知条件确定}
\end{problem}

\begin{answer}
A
\end{answer}

\begin{solution}
设首项为 \(a_1>0\)，公比为 \(q>0\)，且 \(q\neq1\). 则
\[
a_1+a_8-a_4-a_5
=a_1(1+q^7-q^3-q^4)
=a_1(1-q^3)(1-q^4)
\]
当 \(q>1\) 时两个括号均为负，当 \(0<q<1\) 时两个括号均为正，所以乘积始终为正. 因此 \(a_1+a_8>a_4+a_5\)，选 A.
\end{solution}

\begin{problem}
设有如下三个命题：甲：相交两直线 \(l\)，\(m\) 都在平面 \(\alpha\) 内，并且都不在平面 \(\beta\) 内. 乙：\(l\)，\(m\) 之中至少有一条与 \(\beta\) 相交. 丙：\(\alpha\) 与 \(\beta\) 相交. 当甲成立时（\(\quad\)）

\choices
    {乙是丙的充分而不必要的条件}
    {乙是丙的必要而不充分的条件}
    {乙是丙的充分且必要的条件}
    {乙既不是丙的充分条件又不是丙的必要条件}
\end{problem}

\begin{answer}
C
\end{answer}

\begin{solution}
若平面 \(\alpha\) 与 \(\beta\) 相交，交线记为 \(t\). 在平面 \(\alpha\) 内，不在 \(\beta\) 内的直线若不与 \(\beta\) 相交，只能平行于 \(t\). 若 \(l\)，\(m\) 都不与 \(\beta\) 相交，则 \(l\)，\(m\) 都平行于 \(t\)，从而彼此平行，这与 \(l\)，\(m\) 相交矛盾. 所以丙成立时乙成立.

反过来，若 \(l\) 或 \(m\) 与 \(\beta\) 相交，由于该直线在 \(\alpha\) 内，交点同时属于 \(\alpha\) 和 \(\beta\)，故两平面相交，即丙成立. 所以乙与丙互为充分必要条件，选 C.
\end{solution}

\begin{problem}
将数字 \(1\)，\(2\)，\(3\)，\(4\) 填入标号为 \(1\)，\(2\)，\(3\)，\(4\) 的四个方格里，每格填一个数字，则每个方格的标号与所填的数字均不相同的填法有（\(\quad\)）

\choices
    {\(6\) 种}
    {\(9\) 种}
    {\(11\) 种}
    {\(23\) 种}
\end{problem}

\begin{answer}
B
\end{answer}

\begin{solution}
这是四个元素的错排问题. 用容斥原理，先不限制有 \(4!=24\) 种排列；若指定 \(j\) 个数字放入对应方格，剩下数字有 \((4-j)!\) 种排法，因此
\[
4!-\mathrm C_4^1 3!+\mathrm C_4^2 2!-\mathrm C_4^3 1!+\mathrm C_4^4 0!
=24-24+12-4+1=9
\]
故共有 \(9\) 种，选 B.
\end{solution}

\section{填空题}

\begin{problem}
\(\sin\left(\arccos \frac{1}{2}+\arccos \frac{1}{3}\right)=\)\fillinblank{}.
\end{problem}

\begin{answer}
\(\frac{\sqrt{3}+2\sqrt{2}}{6}\)
\end{answer}

\begin{solution}
令 \(\alpha=\arccos\frac12\)，\(\beta=\arccos\frac13\). 由于 \(\alpha\)，\(\beta\in[0,\pi]\)，有
\(\sin\alpha=\sqrt{1-\cos^2\alpha}=\frac{\sqrt3}{2}\)，\(\sin\beta=\sqrt{1-\cos^2\beta}=\frac{2\sqrt2}{3}\).
由正弦的和角公式，
\[
\sin(\alpha+\beta)=\sin\alpha\cos\beta+\cos\alpha\sin\beta
=\frac{\sqrt3}{6}+\frac{\sqrt2}{3}
=\frac{\sqrt3+2\sqrt2}{6}
\]
\end{solution}

\begin{problem}
设 \(a>1\)，则 \(\displaystyle\lim\limits_{n\to\infty}\frac{1-a^{n+1}}{1+a^{n-1}}=\)\fillinblank{}.
\end{problem}

\begin{answer}
\(-a^2\)
\end{answer}

\begin{solution}
因 \(a>1\)，所以 \(a^{n-1}\to+\infty\). 分子、分母同除以 \(a^{n-1}\)，得
\[
\frac{1-a^{n+1}}{1+a^{n-1}}
=\frac{a^{-(n-1)}-a^2}{a^{-(n-1)}+1}
\longrightarrow -a^2
\]
\end{solution}

\begin{problem}
从 \(1\)，\(2\)，\(\cdots\)，\(10\) 这十个数中取出四个数，使它们的和为奇数，共有\fillinblank{} 种取法（用数字作答）.
\end{problem}

\begin{answer}
100
\end{answer}

\begin{solution}
\(1\) 到 \(10\) 中有 \(5\) 个奇数和 \(5\) 个偶数. 四个数的和为奇数，当且仅当选出的奇数个数为 \(1\) 或 \(3\). 因而取法数为
\[
\mathrm C_5^1\mathrm C_5^3+\mathrm C_5^3\mathrm C_5^1=5\cdot10+10\cdot5=100
\]
\end{solution}

\begin{problem}
设 \(f(x)=4^x-2^{x+1}\)，则 \(f^{-1}(0)=\)\fillinblank{}.
\end{problem}

\begin{answer}
1
\end{answer}

\begin{solution}
求 \(f^{-1}(0)\) 即求方程 \(f(x)=0\) 的解. 令 \(t=2^x\)，则 \(t>0\)，且
\[
4^x-2^{x+1}=2^{2x}-2\cdot2^x=t^2-2t=t(t-2)
\]
因此 \(t=0\) 或 \(t=2\)，舍去不可能的 \(t=0\)，得到 \(2^x=2\)，即 \(x=1\). 所以 \(f^{-1}(0)=1\).
\end{solution}

\begin{problem}
建造一个容积为 \(8\mathrm{~m}^3\)，深为 \(2\mathrm{~m}\) 的长方体无盖水池. 如果池底和池壁的造价每平方米分别为 \(120\text{ 元}\) 和 \(80\text{ 元}\)，那么水池的最低总造价为\fillinblank{}\(\text{元}\).
\end{problem}

\begin{answer}
1760
\end{answer}

\begin{solution}
设池底长、宽分别为 \(x\mathrm{~m}\)、\(y\mathrm{~m}\). 由容积和深度得
\(2xy=8\)，\(xy=4\).
池底造价为 \(120xy=480\) 元，四面池壁面积为 \(2\cdot2x+2\cdot2y=4(x+y)\)，池壁造价为 \(320(x+y)\) 元. 因为
\[
x+y\geq2\sqrt{xy}=4
\]
且 \(x=y=2\) 时取等号，所以最低总造价为
\[
480+320\cdot4=1760\text{ 元}
\]
\end{solution}

\begin{problem}
如图，\(ABCD\) 是正方形，\(E\) 是 \(AB\) 的中点，如将 \(\triangle DAE\) 和 \(\triangle CBE\) 分别沿虚线 \(DE\) 和 \(CE\) 折起，使 \(AE\) 与 \(BE\) 重合，记 \(A\) 与 \(B\) 重合的点为 \(P\)，则面 \(PCD\) 与面 \(ECD\) 所成的二面角为\fillinblank{} 度.

\bitmapfigure[width=0.62\linewidth]{img/1993/jing_xiang_e_yun_qiong_gui_liberal/q23_fig1.png}
\end{problem}

\begin{answer}
30
\end{answer}

\begin{solution}
设正方形边长为 \(s\)，令 \(O\) 为 \(CD\) 的中点. 折起后，\(P\) 是 \(A\)，\(B\) 的重合点，所以
\(PD=AD=s\)，\(PC=BC=s\)，\(PE=AE=BE=\frac s2\).
由 \(PC=PD\)，得 \(PO\perp CD\)，且
\[
PO=\sqrt{PD^2-OD^2}=\sqrt{s^2-\left(\frac s2\right)^2}=\frac{\sqrt3}{2}s
\]
在原正方形的平面中，\(OE=s\)，且 \(OE\perp CD\). 因此 \(\angle POE\) 是两个平面 \(PCD\)、\(ECD\) 所成的二面角的平面角. 在三角形 \(POE\) 中，
\[
PE^2=PO^2+OE^2-2\,PO\,OE\cos\angle POE
\]
代入 \(PE=s/2\)、\(PO=\sqrt3s/2\)、\(OE=s\)，得
\[
\frac{s^2}{4}=\frac{3s^2}{4}+s^2-\sqrt3s^2\cos\angle POE
\]
所以 \(\cos\angle POE=\sqrt3/2\)，从而二面角为 \(30^\circ\).
\end{solution}

\section{解答题}

\begin{problem}
求 \(\tan 20^{\circ}+4\sin 20^{\circ}\) 的值.
\end{problem}

\begin{solution}
令 \(s=\sin20^\circ\)，\(c=\cos20^\circ\). 由三倍角公式，
\(2s(4c^2-1)=2\sin60^\circ=\sqrt3\)，\(8c^3-6c=2\cos60^\circ=1\).
于是 \(2c(4c^2-1)=8c^3-2c=1+4c\). 因此
\[
\tan20^\circ+4\sin20^\circ
=\frac{s}{c}+4s
=\frac{s(1+4c)}{c}
=2s(4c^2-1)=\sqrt3
\]
\end{solution}

\begin{problem}
已知 \(f(x)=\log_a \frac{1+x}{1-x}\)（\(a>0\)，\(a\neq 1\)）.

\begin{enumerate}
\item 求 \(f(x)\) 的定义域；
\item 判断 \(f(x)\) 的奇偶性并予以证明；
\item 求使 \(f(x)>0\) 的 \(x\) 取值范围.
\end{enumerate}
\end{problem}

\begin{solution}
\begin{enumerate}
\item 对数的真数须为正，即 \(\dfrac{1+x}{1-x}>0\). 分子、分母的零点分别为 \(-1\)、\(1\)，作符号判断可得 \(-1<x<1\)，所以定义域为 \((-1,1)\).
\item 定义域 \((-1,1)\) 关于原点对称，且对任意 \(x\in(-1,1)\)，有
\[
f(-x)=\log_a\frac{1-x}{1+x}=-\log_a\frac{1+x}{1-x}=-f(x)
\]
所以 \(f(x)\) 是奇函数.
\item 当 \(a>1\) 时，\(f(x)>0\) 等价于 \(\dfrac{1+x}{1-x}>1\)，解得 \(0<x<1\)；当 \(0<a<1\) 时，不等号方向相反，解得 \(-1<x<0\).
\end{enumerate}
\end{solution}

\begin{problem}
已知数列 \(\frac{8\cdot 1}{1^2\cdot 3^2}\)，\(\frac{8\cdot 2}{3^2\cdot 5^2}\)，\(\cdots\)，\(\frac{8n}{(2n-1)^2(2n+1)^2}\)，\(\cdots\). \(S_n\) 为其前 \(n\) 项和. 计算得 \(S_1=\frac{8}{9}\)，\(S_2=\frac{24}{25}\)，\(S_3=\frac{48}{49}\)，\(S_4=\frac{80}{81}\). 观察上述结果，推测出计算 \(S_n\) 的公式，并用数学归纳法加以证明.
\end{problem}

\begin{solution}
由 \(S_1\)，\(S_2\)，\(S_3\)，\(S_4\) 的形式，猜想
\[
S_n=\frac{4n(n+1)}{(2n+1)^2}
\]
当 \(n=1\) 时，原数列首项为 \(S_1=8/9\)，而
\[
\frac{4\cdot1\cdot2}{(2\cdot1+1)^2}=\frac89
\]
所以公式成立.

假设当 \(n=k\) 时公式成立，即
\[
S_k=\frac{4k(k+1)}{(2k+1)^2}
\]
由归纳假设 \(S_k=\dfrac{4k(k+1)}{(2k+1)^2}\)，并注意到
\[
k(2k+3)^2+2=(k+2)(2k+1)^2
\]
因为等式两边展开后均为 \(4k^3+12k^2+9k+2\). 因此
\[
\begin{gathered}
S_{k+1}=S_k+\frac{8(k+1)}{(2k+1)^2(2k+3)^2}\\
=\frac{4k(k+1)}{(2k+1)^2}
+\frac{8(k+1)}{(2k+1)^2(2k+3)^2}\\
=\frac{4(k+1)\left[k(2k+3)^2+2\right]}
{(2k+1)^2(2k+3)^2}\\
=\frac{4(k+1)(k+2)}{(2k+3)^2}.
\end{gathered}
\]
即当 \(n=k+1\) 时公式也成立. 根据数学归纳法，对任意正整数 \(n\)，都有
\[
S_n=\frac{4n(n+1)}{(2n+1)^2}
\]
\end{solution}

\begin{problem}
已知：平面 \(\alpha\cap\beta=\) 直线 \(a\). \(\alpha\)，\(\beta\) 同垂直于平面 \(\gamma\)，又同平行于直线 \(b\). 求证：

\begin{enumerate}
\item \(a\perp\gamma\)；
\item \(b\perp\gamma\).
\end{enumerate}
\end{problem}

\begin{solution}
\begin{enumerate}
\item 设 \(n_\alpha\)，\(n_\beta\)，\(n_\gamma\) 分别为平面 \(\alpha\)，\(\beta\)，\(\gamma\) 的法向量. 因为 \(\alpha\perp\gamma\)，所以 \(n_\alpha\perp n_\gamma\)，即 \(n_\alpha\) 的方向在平面 \(\gamma\) 内；同理 \(n_\beta\) 的方向也在平面 \(\gamma\) 内. 又因 \(\alpha\) 与 \(\beta\) 相交，\(n_\alpha\)，\(n_\beta\) 不平行；而 \(a=\alpha\cap\beta\)，故 \(a\) 的方向同时垂直于 \(n_\alpha\)，\(n_\beta\)，从而平行于 \(n_\alpha\times n_\beta\). 两个不平行且都在平面 \(\gamma\) 内的向量的叉积垂直于平面 \(\gamma\)，所以 \(a\perp\gamma\).
\item 直线 \(b\) 同时平行于平面 \(\alpha\)，\(\beta\)，所以 \(b\) 的方向向量同时垂直于 \(n_\alpha\)，\(n_\beta\). 因而 \(b\) 的方向平行于 \(n_\alpha\times n_\beta\)，即 \(b\parallel a\). 由第（1）问 \(a\perp\gamma\)，所以 \(b\perp\gamma\).
\end{enumerate}
\end{solution}

\begin{problem}
在面积为 \(1\) 的 \(\triangle PMN\) 中，\(\tan \angle PMN=\frac{1}{2}\)，\(\tan \angle MNP=-2\). 建立适当的坐标系，求以 \(M\)，\(N\) 为焦点且过点 \(P\) 的椭圆方程.
\end{problem}

\begin{solution}
建立直角坐标系，使 \(M\)，\(N\) 分别为 \(x\) 轴上的两点，且 \(M=(0,0)\)，\(N=(d,0)\). 取 \(P\) 在 \(MN\) 的上方，则由 \(\tan\angle PMN=\frac12\)，直线 \(MP\) 的斜率为 \(\frac12\)；又因 \(\tan\angle MNP=-2\)，角 \(\angle MNP\) 为钝角，直线 \(NP\) 的斜率为 \(2\). 故点 \(P\) 的坐标为 \(\left(\frac{4d}{3},\frac{2d}{3}\right)\). 由三角形面积为 \(1\)，得 \(\frac12d\cdot\frac{2d}{3}=1\)，所以 \(d=\sqrt{3}\).

于是 \(MP=\dfrac{2\sqrt{15}}{3}\)，\(NP=\dfrac{\sqrt{15}}{3}\)，从而椭圆的长轴长为 \(MP+NP=\sqrt{15}\)，半长轴 \(a=\dfrac{\sqrt{15}}{2}\)，半焦距 \(c=\dfrac{MN}{2}=\dfrac{\sqrt{3}}{2}\)，半短轴满足 \(b^2=a^2-c^2=3\). 若取椭圆中心为原点，则 \(M\)，\(N\) 的坐标为 \(\left(-\frac{\sqrt{3}}{2},0\right)\)，\(\left(\frac{\sqrt{3}}{2},0\right)\)，故所求椭圆方程为
\[
\frac{x^2}{15/4}+\frac{y^2}{3}=1
\]
\end{solution}
